\bookchapter{Κεφάλαιο 18: Αντίγραφα ασφαλείας δεδο$\mu$ένων και $\mu$εταφορά}

\chapterintro{Μια γειτόνισσά $\mu$ου, το κινητό της έπεσε στην πισίνα και δεν ανοίγει. Δεν την πείραξε τόσο το ίδιο το κινητό, αλλά οι φωτογραφίες τριών ετών χωρίς backup — όλα τα στάδια του παιδιού της, από τη γέννηση $\mu$έχρι τα πρώτα του βή$\mu$ατα. Όταν $\mu$ου το είπε, ένιωσα άσχη$\mu$α, γιατί αυτό θα $\mu$πορούσε να είχε αποφευχθεί. Σε αυτή την ενότητα θα ασφαλίσου$\mu$ε σοβαρά τα δεδο$\mu$ένα $\mu$ας.}

\booksection{I. Πρώτα ξεκαθαρίστε: τι χάνεται αν χαθεί}

Διαφορετικά δεδο$\mu$ένα έχουν διαφορετική “κατάσταση επιβίωσης”:

\begin{itemize}
  \item \textbf{αποθηκεύονται τοπικά στο τηλέφωνο — αν χαθεί, χάνονται οριστικά}. Φωτογραφίες/βίντεο (χωρίς sync), notes, SMS, κλήσεις, λήψεις.
  \item \textbf{Αποθηκευ$\mu$ένα στον server της App, τα ανακτάτε $\mu$ε νέα σύνδεση στο κινητό}: WeChat φίλοι, Taobao παραγγελίες, τραπεζικές συναλλαγές — σε servers, όχι στο κινητό.
  \item \textbf{Σε γκρίζα ζώνη, απαιτεί ειδικές ενέργειες}: WeChat ιστορικό (σε server αλλά όχι download), τραπεζικά device credentials, 2FA App secrets.
\end{itemize}

Επο$\mu$ένως αυτές που πρέπει να προστατεύσετε είναι η πρώτη και η τρίτη κατηγορία.

\booksection{II. Πλήρες αντίγραφο συσκευής: όλα $\mu$ε $\mu$ια φορά}

Το πλήρες αντίγραφο είναι σαν “snapshot” της συσκευής — ρυθ$\mu$ίσεις, διάταξη, App, φωτογραφίες, επαφές, όλα $\mu$αζί.

\booksubsection{Επιλογή 1: Τοπικό αντίγραφο σε υπολογιστή (πιο αξιόπιστο, για όσους έχουν υπολογιστή)}

\begin{itemize}
  \item \textbf{iPhone}. Καλώδιο σε υπολογιστή, iTunes/Finder, "Back Up Now". Τσεκάρετε "Encrypt", αλλιώς κωδικοί/υγεία δεν $\mu$παίνουν.
  \item \textbf{Android}: Κάθε $\mu$άρκα έχει PC software: Huawei Assistant, Xiaomi Assistant, Samsung Smart Switch. Συνδέστε και ακολουθήστε.
  \item Το τοπικό αντίγραφο γλιτώνει χρή$\mu$ατα για cloud και η ανάκτηση είναι πολύ πιο γρήγορη. Μειονέκτη$\mu$α: πρέπει να θυ$\mu$άστε να συνδέετε τακτικά — οι περισσότεροι δεν το κάνουν.
\end{itemize}

\booksubsection{Επιλογή 2: Cloud αντίγραφο (πιο ξεκούραστο, για όλους)}

\begin{itemize}
  \item \textbf{iPhone}. Ρυθ$\mu$ίσεις όνο$\mu$α iCloud Αντίγραφο On. Φορτιστή, Wi-Fi, lock = αυτό$\mu$ατο. 5GB δεν αρκούν, 50GB $\mu$ε 6 γουάν.
  \item \textbf{Android}: Κάθε $\mu$άρκα έχει cloud. Huawei "cloud backup", Xiaomi "Xiaomi Cloud", OPPO/vivo παρό$\mu$οια. Και Google backup (Ρυθ$\mu$ίσεις Σύστη$\mu$α Backup), ανοίξτε το.
  \item Το πλεονέκτη$\mu$α του cloud backup είναι "χωρίς να σας νοιάζει", γίνεται αυτό$\mu$ατα. Το $\mu$ειονέκτη$\mu$α: η ανάκτηση θέλει πολλή λήψη και ο δωρεάν χώρος δεν αρκεί.
\end{itemize}

\booksection{III. Στρατηγική κατηγοριοποιη$\mu$ένου αντιγράφου}

Το πλήρες αντίγραφο είναι δίχτυ ασφαλείας, αλλά $\mu$ερικά πράγ$\mu$ατα αξίζουν ξεχωριστό αντίγραφο:

\booksubsection{Φωτογραφίες και βίντεο (το πιο πολύτι$\mu$ο ψηφιακό περιουσιακό)}

\begin{itemize}
  \item \textbf{Σχέδιο διπλής ασφάλειας}: Sync iCloud/Google Photos, plus κάθε λίγους $\mu$ήνες manual copy σε υπολογιστή/δίσκο. Cloud + τοπικό.
  \item \textbf{Λύση χωρίς κόστος}: Google Photos 15GB δωρεάν, ρυθ$\mu$ίστε "compressed" για χώρο. Ή κάθε $\mu$ήνα $\mu$ε καλώδιο σε υπολογιστή.
\end{itemize}

\booksubsection{Επαφές (ευκολόθεα αλλά ζωτικής ση$\mu$ασίας)}

\begin{itemize}
  \item Οι επαφές iPhone συγχρονίζονται από προεπιλογή $\mu$ε iCloud, αρκεί ο συγχρονισ$\mu$ός να είναι ανοιχτός. Για περισσότερη ασφάλεια, εξάγετε vCard από το iCloud.com στον υπολογιστή.
  \item Οι επαφές Android συνήθως συγχρονίζονται $\mu$ε Google ή το cloud της επωνυ$\mu$ίας. Πηγαίνετε στις ρυθ$\mu$ίσεις της εφαρ$\mu$ογής Επαφών και εξάγετε αρχείο .vcf, είναι ένα πλήρες αντίγραφο ασφαλείας.
\end{itemize}

\booksubsection{Ιστορικό WeChat (ξεχωριστή διαχείριση, υψίστης ση$\mu$ασίας)}

Αυτό χρειάζεται ξεχωριστή αναφορά, γιατί στο πλήρες αντίγραφο το ιστορικό WeChat είναι συχνά ελλιπές και τα εργαλεία $\mu$εταφοράς δεν το κινούν.

\begin{itemize}
  \item \textbf{Μέθοδος 1: Μεταφορά σε νέο κινητό}. WeChat Εγώ Ρυθ$\mu$ίσεις Γενικά Μεταφορά \& αντίγραφο ιστορικού Μεταφορά σε τηλέφωνο/tablet. Νέο και παλιό στο ίδιο Wi-Fi, σκανάρισ$\mu$α. Μην αποσυνδέσετε το WeChat στο παλιό $\mu$έχρι να ολοκληρωθεί.
  \item \textbf{Μέθοδος 2: Αντίγραφο σε υπολογιστή}. Εγκαταστήστε WeChat PC, κάτω αριστερά τρεις γρα$\mu$$\mu$ές Μεταφορά \& αντίγραφο Αντίγραφο \& ανάκτηση Αντίγραφο ιστορικού σε υπολογιστή. Αργότερα ανακτάται στο νέο κινητό.
  \item \textbf{Μέθοδος 3: Αντίγραφο στο cloud}. Το WeChat υποστηρίζει αντίγραφο στο iCloud ή στο cloud του κινητού, αλλά: ανακτάται $\mu$όνο στο τρέχον κινητό. Για αλλαγή κινητού χρησι$\mu$οποιήστε τις δύο πρώτες $\mu$εθόδους.
\end{itemize}

\booksection{IV. Οδηγός αλλαγής συσκευής: από παλιό σε νέο τηλέφωνο}

Η αλλαγή κινητού είναι χαρά, αλλά αν δεν $\mu$εταφερθούν τα δεδο$\mu$ένα γίνεται πονοκέφαλος. Τρεις περιπτώσεις:

\booksubsection{Αναβάθ$\mu$ιση ίδιας επωνυ$\mu$ίας (πιο απλή)}

\begin{itemize}
  \item \textbf{iPhone νέο iPhone}: Ενεργοποίηση $\mu$ε παλιό δίπλα, "Quick Start", scan. Όλα $\mu$εταφέρονται, και κωδικοί Wi-Fi. 20-30 λεπτά.
  \item \textbf{Android ίδιας επωνυ$\mu$ίας}. Εργαλεία $\mu$άρκας: Huawei "Phone Clone", Xiaomi "Xiaomi Transfer", OPPO "Phone Move", vivo "EasyShare". Ίδια βή$\mu$ατα: App, scan, select, send.
\end{itemize}

\booksubsection{Αλλαγή $\mu$άρκας (Android σε Android)}

Με εργαλεία τρίτων ή App παλιάς $\mu$άρκας (οι περισσότερες $\mu$εταφέρουν σε άλλο Android). Μεταφέρονται: επαφές, φωτογραφίες, SMS, κλήσεις, $\mu$ερικές App. Δεν $\mu$εταφέρονται: login στις App, ιστορικό WeChat (χρησι$\mu$οποιήστε τη $\mu$εταφορά του WeChat).

\booksubsection{Αλλαγή πλατφόρ$\mu$ας (Android σε iPhone ή αντίστροφα)}

\begin{itemize}
  \item \textbf{Android iPhone}. Apple "Move to iOS" App, στο Android. Μεταφέρει επαφές, φωτογραφίες, SMS, όχι App data. WeChat ξεχωριστά.
  \item \textbf{iPhone Android}: Κάθε $\mu$άρκα έχει εργαλεία, αλλά cross-platform χειρότερο. Πρώτα: επαφές, φωτογραφίες, WeChat. Υπόλοιπα αργότερα.
\end{itemize}

\booksection{V. Χρυσός κανόνας αντιγράφων: κανόνας 3-2-1}

Κλασική αρχή ασφάλειας δεδο$\mu$ένων, εύκολη:

\begin{itemize}
  \item \textbf{3}. Ση$\mu$αντικά δεδο$\mu$ένα 3 αντίγραφα (1 original + 2).
  \item \textbf{2}: Αντίγραφα σε δύο $\mu$έσα (cloud + εξωτερικό δίσκο).
  \item \textbf{1}. Ένα αντίγραφο σε άλλο $\mu$έρος (φωτιά, κλοπή).
\end{itemize}

Για τους απλούς ανθρώπους, να έχετε “ένα στο cloud, ένα στον υπολογιστή” είναι ήδη πολύ καλό. Μην το κάνετε πολύ περίπλοκο, το κλειδί είναι να το κάνετε πραγ$\mu$ατικά.

\begin{tcbwarn}{Ειδική προσοχή στο ιστορικό συνο$\mu$ιλιών WeChat}
Το ιστορικό συνο$\mu$ιλιών WeChat δεν $\mu$εταφέρεται $\mu$έσω κανενός εργαλείου $\mu$εταφοράς κινητού! Δεν καλύπτεται από τα πλήρη αντίγραφα ασφαλείας. Πρέπει να χρησι$\mu$οποιήσετε τη λειτουργία «Μεταφορά ιστορικού συνο$\mu$ιλιών» του WeChat, ή αντίγραφο στην έκδοση PC. Πολλοί ανακαλύπτουν ότι το ιστορικό χάθηκε $\mu$ετά την αλλαγή κινητού, τότε είναι πολύ αργά — το παλιό $\mu$πορεί να έχει ήδη επανέλθει στις εργοστασιακές ρυθ$\mu$ίσεις.

\end{tcbwarn}

\begin{tcbpractice}{Άσκηση: $\mu$ια πλήρης προπόνηση αντιγράφου}
\begin{enumerate}
  \item Ελέγξτε αν iCloud ή cloud αντιγραφή είναι ανοιχτά, δείτε τελευταία φορά.
  \item Χειροκίνητη δη$\mu$ιουργία αντιγράφου φωτογραφιών: συνδέστε το κινητό $\mu$ε καλώδιο δεδο$\mu$ένων στον υπολογιστή, αντιγράψτε τους φακέλους DCIM και Pictures στον υπολογιστή.
  \item Δη$\mu$ιουργία αντιγράφου επαφών: χρήστες iPhone πηγαίνετε στο iCloud.com εξάγετε vCard· χρήστες Android στις ρυθ$\mu$ίσεις της εφαρ$\mu$ογής Επαφών εξάγετε αρχείο .vcf.
  \item Δη$\mu$ιουργία αντιγράφου ασφαλείας ιστορικού WeChat: ανοίξτε το WeChat Εγώ Ρυθ$\mu$ίσεις Γενικά Μεταφορά και δη$\mu$ιουργία αντιγράφου ιστορικού συνο$\mu$ιλιών Δη$\mu$ιουργία αντιγράφου στον υπολογιστή.
  \item Αντιγράψτε τα εξαγό$\mu$ενα vCard και φωτογραφίες σε εξωτερικό δίσκο ή USB, για "ένα αντίγραφο εκτός δικτύου".
\end{enumerate}

\end{tcbpractice}
